\section{Conclusion}
\label{sec:conclusion}

We tested four specific claims about how pre-training mixture proportions shape performance, using 81 small language models trained on 53 controlled mixtures of four domains.
Held-out loss on a domain follows a log-linear law in that domain's share, with no saturation at 20\%, and a law fitted on 29 mixtures predicts loss on unseen mixtures.
Downstream accuracy does not follow the hypothesized laws.
It is flat within noise across interior mixtures, code share does not correlate with reasoning scores ($r \approx 0$, not $> 0.7$), and web text above 50\% does not measurably hurt knowledge tasks.
The one robust downstream effect is that a corpus without natural language is clearly worse on knowledge tasks.

The key takeaway is that, at this scale, a law for loss is not a law for downstream performance.

Several steps follow.
The sweep design should be repeated at 100--500M parameters and at least 1B tokens, where zero-shot benchmarks carry signal; the coarse grid $\{0, 0.05, 0.2, 0.5, 1\}$ with three seeds was the informative part of our design.
Future runs should use at least five seeds per point and continuous metrics such as log-likelihood margins, so that 0.5-point effects are detectable.
Varying the token budget would separate share from the amount of in-domain data and test whether the log-share slopes are budget-invariant.
The region between 0\% and 2\%, where roughly half of each domain's loss benefit is realized, deserves a finer grid.
Open questions remain: whether the flatness persists once models can use domain knowledge zero-shot, whether the \listops--code association is a bracket-syntax effect, and why pre-training makes arithmetic harder to learn for some seeds.
